\documentclass[11pt, a4paper]{article}

% --- Packages ---
\usepackage[margin=1in]{geometry}    % Standard 1-inch margins
\usepackage{amsmath, amssymb}        % Math typesetting
\usepackage{graphicx}                % Image handling
\usepackage{hyperref}                % Clickable links
\usepackage{xcolor}                  % Color support
\usepackage{booktabs}                % Professional tables
\usepackage{cite}                    % Citation management
\usepackage{minted}
\usepackage{listings}
% --- Link Setup ---
\hypersetup{
    colorlinks=true,
    linkcolor=blue,
    filecolor=magenta,      
    urlcolor=cyan,
    citecolor=blue,
}

% --- Title Information ---
\title{\textbf{Project 1 Report}\\ \large INFSCI 2591 - Algorithm Design}
\author{Abhay Kumara Sri Krishna Nandiraju \\ University of Pittsburgh}
\date{Spring 2026}

\begin{document}

\maketitle

% \begin{abstract}
% This is a brief summary of your project. State the problem you are solving, the methodology or architecture you used, and the primary results or performance metrics you achieved. Keep it concise, usually around 150-200 words.
% \end{abstract}

% --- Main Content ---
\section{Merging Sorted Subarrays}
Design and implement your own algorithm that takes the array A with size m+n as input where:
Subarray A[1], A[2],...A[m] sorted in ascending order and 
Subarray A[m+1], A[m+2],...A[n] sorted in ascending order
and merges the two subarrays using an auxiliary array Aux of size $min\{m,n\}$ back into array A sorted in ascending order. You must design and implement your own sorting function. Use of sorting functions in libraries is not permitted. 
\begin{verbatim}
Test Case 1: {} and {3, 7, 9}
Test Case 2: {2, 7, 9} and {1}
Test Case 3: {1, 7, 10, 15} and {3, 8, 12, 18}
Test Case 4: {1, 3, 5, 5, 15, 18, 21} and {5, 5, 6, 8, 10, 12, 16, 17, 17, 20, 25, 28}
\end{verbatim}


% \subsection{Code}
% \begin{minted}{python}
% def merge_optimized(A, m, n):
%     # if either subarray size is 0 then we already the sorted array A
%     if m == 0 or n == 0:
%         return A

%     if m <= n:
%         # if left subarray size(m) < right subarray size(n)
%         # auxiliary array size = m
%         aux = A[:m]
%         i, j, k = 0, m, 0
%         while i < m and j < m + n:
%             if aux[i] <= A[j]:
%                 A[k] = aux[i]
%                 i += 1
%             else:
%                 A[k] = A[j]
%                 j += 1
%             k += 1
%         # i am copying the remaining elements from auxiliary
%         while i < m:
%             A[k] = aux[i]
%             i += 1
%             k += 1

%     else:
%         # if left subarray size(m) > right subarray size(n)
%         # auxiliary array size =n
%         aux = A[m : m + n]
%         i, j, k = n - 1, m - 1, m + n - 1
%         while i >= 0 and j >= 0:
%             if aux[i] >= A[j]:
%                 A[k] = aux[i]
%                 i -= 1
%             else:
%                 A[k] = A[j]
%                 j -= 1
%             k -= 1
%         # copying remaining elements from auxiliary
%         while i >= 0:
%             A[k] = aux[i]
%             i -= 1
%             k -= 1

%     return A


% test_cases_p1 = [
%     (0, 3, [3,7,9]),
%     (3, 1, [2, 7, 9, 1]),
%     (4, 4, [1, 7, 10, 15, 3, 8, 12, 18]),
%     (7,12, [1, 3, 5, 5, 15, 18, 21, 5, 5, 6, 8, 10, 12, 16, 17, 17, 20, 25, 28])
% ]

% print("--- Problem 1 Results ---")
% for idx, (m, n, A) in enumerate(test_cases_p1):
%     original = list(A)
%     merged = merge_optimized(A, m, n)
%     print(f"Test Case {idx+1}: \n  Original: {original}\n  Merged:   {merged}\n")
    
% \end{minted}

Code file \texttt{problem1.py} contains the implementation of the above problem. The results from running the code gives sorted array A as follows:
\begin{verbatim}
python problem1.py 
--- Problem 1 Results ---
Test Case 1: 
  Original: [3, 7, 9]
  Merged:   [3, 7, 9]

Test Case 2: 
  Original: [2, 7, 9, 1]
  Merged:   [1, 2, 7, 9]

Test Case 3: 
  Original: [1, 7, 10, 15, 3, 8, 12, 18]
  Merged:   [1, 3, 7, 8, 10, 12, 15, 18]

Test Case 4: 
  Original: [1, 3, 5, 5, 15, 18, 21, 5, 5, 6, 8, 10, 12, 16, 17, 17, 20, 25, 28]
  Merged:   [1, 3, 5, 5, 5, 5, 6, 8, 10, 12, 15, 16, 17, 17, 18, 20, 21, 25, 28]
\end{verbatim}


\newpage
\section{Multiplication}
Design and implement your own algorithms, one for Ala Carte Multiplication ( to break down the multiplication process into smaller, more manageable parts) and one for Rectangle Multiplication (to create a grid (or rectangle) where the area represents the product of the two numbers). Your algorithms must allow for both positive and negative multiplicands and multipliers.
\begin{verbatim}
Test Case 1: 7000 * 7294
Test Case 2: 25 * 5038385
Test Case 3: -59724 * 783
Test Case 4: 8516 * -82147953548159344
Test Case 5: 45952456856498465985 * 98654651986546519856
Test Case 6: -45952456856498465985 * -98654651986546519856
\end{verbatim}
\subsection{AlaCarte Multiplication}
Code file \texttt{problem2\_alacarte.py} contains the implementation of Ala Carte Multiplication. The results from running the code gives the products as follows:
\begin{verbatim}
python problem2_alacarte.py 
--- Problem 2 Alacarte Results ---
Test Case 1: 
  Num1: 7000
  Num2: 7294
  Product: 51058000

Test Case 2: 
  Num1: 25
  Num2: 5038385
  Product: 125959625

Test Case 3: 
  Num1: -59724
  Num2: 783
  Product: -46763892

Test Case 4: 
  Num1: 8516
  Num2: -82147953548159344
  Product: -699571972416124973504

Test Case 5: 
  Num1: 45952456856498465985
  Num2: 98654651986546519856
  Product: 4533423639104649634397093450504343098160

Test Case 6: 
  Num1: -45952456856498465985
  Num2: -98654651986546519856
  Product: 4533423639104649634397093450504343098160
\end{verbatim}

\subsection{Rectangle Multiplication}
Code file \texttt{problem2\_rectangle.py} contains the implementation of Rectangle Multiplication. The results from running the code gives the products as follows:
\begin{verbatim}
python problem2_rectangle.py 
--- Problem 2 Rectangle Results ---
Test Case 1: 
  Num1: 7000
  Num2: 7294
  Product: 51058000

Test Case 2: 
  Num1: 25
  Num2: 5038385
  Product: 125959625

Test Case 3: 
  Num1: -59724
  Num2: 783
  Product: -46763892

Test Case 4: 
  Num1: 8516
  Num2: -82147953548159344
  Product: -699571972416124973504

Test Case 5: 
  Num1: 45952456856498465985
  Num2: 98654651986546519856
  Product: 4533423639104649634397093450504343098160

Test Case 6: 
  Num1: -45952456856498465985
  Num2: -98654651986546519856
  Product: 4533423639104649634397093450504343098160

\end{verbatim}

\textbf{Note:} In Python, the built-in integer type can handle arbitrarily
large integers, so we can directly compute the products without worrying about overflow.
If the algorithms are implemented in C++ or Java, we might need to handle large integers during computation.
\end{document}